% =====================================================================
%  GMN2 §6  WKB 三角剖分（The WKB triangulation）
%  原文：wkb-paper.tex 4708–5109 行；印刷 p.76–83；编号式 (6.1)–(6.10) 共 10 条
%        （原文 (6.6)(6.7) 同出一个 align，按施工指令拆成两个 equation）；脚注 3 条；图 23–29 共 7 幅
%  块内引用用 \eqref / \ref；跨块引用硬编码原编号并加 % TODO
% =====================================================================

\section{WKB 三角剖分}
\label{sec:WKB-triangulation}

上一节我们回顾了 Fock–Goncharov 坐标 $X^T_E$ 的一些基本性质。那么，怎么把它们与我们想要的坐标 $X^\vartheta_\gamma$ 联系起来呢？一旦试图把两者联系起来，一个问题立刻浮现：$X^T_E$ 依赖于三角剖分 $T$，而 $X^\vartheta_\gamma$ 不依赖任何三角剖分，转依赖角 $\vartheta$。所以若想把两者等同起来，就需要一种办法，选出一个依赖于 $\vartheta$ 的典范三角剖分 $T$。这个选择不能任意作：归根到底，我们希望把 $X^\vartheta_\gamma$ 设计得满足第 2 节列出的那些性质。% TODO 合稿改 \ref{sec:review}
特别地，我们想控制 $X^\vartheta_\gamma$ 在 $\zeta \to 0$ 极限下的渐近行为。

这些渐近行为怎么确定呢？基本想法非常简单。当 $\zeta\to0$ 时，$\CA$ 由它的领头项 $R\varphi/\zeta$ 主导。在 $C$ 的任意一点，$\varphi$ 有两个本征值 $\pm\lambda$，其中 $\lambda$ 是 $C$ 上的多值一形式，或是谱曲线 $\Sigma$ 上的单值一形式。「WKB 近似」大致是说：在使
\begin{equation}
\varphi = \begin{pmatrix} \lambda & 0 \\ 0 & -\lambda \end{pmatrix},
\end{equation}
成立的那个规范中，存在两个独立的、形如
\begin{equation} \label{eq:wkb-rough}
\psi^{(1)} \sim \begin{pmatrix} e^{- \frac{R}{\zeta} \int^z \lambda } \\ 0 \end{pmatrix}, \quad \psi^{(2)} \sim \begin{pmatrix} 0 \\ e^{\frac{R}{\zeta} \int^z \lambda } \end{pmatrix}.
\end{equation}
的近似 $\CA$-平坦截面。于是人们也许会预期：$\zeta\to0$ 极限下平行移动的计算，将归结为计算一形式 $\lambda$ 在 $\Sigma$ 上的周期。

把这个想法仔细做出来，就会发现并非所有三角剖分生来平等。的确，固定二次微分 $\lambda^2$ 并选定一个角 $\vartheta$，就存在唯一的「WKB 三角剖分」$T_\WKB(\vartheta, \lambda^2)$，使得当 $\zeta$ 位于以 $e^{i\vartheta}$ 为中心的那个半平面 $\IH_\vartheta$ 内时，WKB 近似给出正确的渐近行为。本节我们定义这个三角剖分，并描述它的一些基本性质。

本节通篇假定 $\lambda^2$ 固定且一般，即只有单零点。

\subsection{WKB 曲线}

固定任意 $\vartheta \in \IR/2\pi\IZ$。我们定义一条\emph{角为 $\vartheta$ 的 WKB 曲线}为 $C$ 中一条曲线，其切向量 $\partial_t$ 满足：沿曲线处处有
\begin{equation} \label{eq:wkb-curve}
\lambda \cdot \partial_t \in e^{i\vartheta} \IR^\times .
\end{equation}
（$\lambda$ 在 $C$ 上只定义到相差一个符号，但这个含糊之处对此定义无关紧要。）这些曲线定义了\emph{角为 $\vartheta$ 的 WKB 叶状结构}。

（注意 \eqref{eq:wkb-curve} 中的参数 $t$，与第 3 节和第 9 节所用的坐标 $t$ 毫无关系。% TODO 合稿改 \ref{sec:physics} 与 \ref{section:Superconformal}
此外，当我们谈论截面沿 WKB 曲线的「指数式增长」时，所取的是使 $\lambda \cdot \partial_t = \pm e^{i\vartheta}$ 的参数化。）

\subsection{WKB 曲线的局部行为} \label{sec:wkb-local}

在 $C$ 的一般点附近，WKB 叶状结构的局部行为是平凡的：用局部坐标 $w = \int \lambda$ 来说，它不过是由直线 $\Im e^{-i\vartheta} w = const$ 给出的叶状结构。

现在考虑某个奇点 $\CP$ 附近的行为。我们取坐标使该点位于 $z=0$，并在 (4.1) 中取 $\rho = 2m\sigma^3$，$m \in \IC$。% TODO 合稿改 \eqref{eq:reg-sing-bc-1}
固定 $\lambda$ 的一个分支，则在 $z=0$ 附近近似地有 $\lambda = \frac{m\,dz}{z}$。

令
\begin{equation} \label{eq:defeps}
\varepsilon(\vartheta) = \sgn \left( \re e^{-i\vartheta} m \right)
\end{equation}
（约定 $\sgn 0 = 0$）。若 $\varepsilon(\vartheta) = 0$，则没有 WKB 曲线进入该奇点（此时附近的 WKB 曲线是绕着它转的圆）；我们暂且假定不处于这个退化情形，第 7.6.3 节会回到这一点。% TODO 合稿改 \ref{sec:no-flavor-jump}
只要 $\varepsilon(\vartheta) \neq 0$，最一般的 WKB 曲线就是一条对数螺线，可以把它参数化成
\begin{equation} \label{eq:log-spiral}
z(t) = z_0\, e^{- \varepsilon(\vartheta) \frac{e^{i\vartheta}}{m} t}.
\end{equation}
这条曲线当 $t \to \infty$ 时进入奇点。注意 $-\varepsilon(\vartheta)$ 就是 $e^{-i\vartheta}\lambda \cdot \partial_t$ 的符号。特别地，对进入同一奇点的所有 WKB 曲线，这个符号都相同。

\insfig{5in}{foliation-pole}{WKB 曲线在奇点附近的行为。}

我们还会关心 WKB 叶状结构在 $\lambda^2$ 的单零点附近的行为，这种点也称为\emph{转向点}。在这样的点上，叶状结构变成奇异的。有三条分离 WKB 曲线从转向点发出，如图 23 所示。% TODO 合稿改 \ref{fig:foliation-zero}
\insfig{2.05in}{foliation-zero}{WKB 叶状结构在转向点附近的行为。一般 WKB 曲线画成细黑线，分离 WKB 曲线画成较粗的红线。}
这些曲线我们下面还会进一步讨论。

\subsection{WKB 曲线的整体行为}\label{subsec:GlobalBehaviorWKB}

现在来讨论 WKB 叶状结构的整体行为。对黎曼面上一般的叶状结构，叶的行为可以相当狂野。幸运的是，WKB 叶状结构要容易控制得多。我们把 WKB 曲线作如下分类：
\begin{itemize}
\item 一条\emph{一般} WKB 曲线，沿两个方向都渐近于奇点（两个方向可以是同一个奇点）。
\item 一条\emph{分离} WKB 曲线，沿一个方向渐近于转向点，沿另一个方向渐近于奇点。
\item 一条\emph{有限} WKB 曲线，沿两个方向都渐近于转向点（可以是同一个转向点），或者是闭合的。
\item 一条\emph{发散} WKB 曲线，不闭合，且沿一个或两个方向不趋于任何极限。
\end{itemize}
在下文中，存在有限 WKB 曲线的那些 $\vartheta$ 值将起非常特殊的作用。它们对应 BPS 态出现、WKB 三角剖分跳变的位置。所以暂时假定不存在有限 WKB 曲线。此时，正如我们下面要证明的，也不存在发散曲线。

假设（为导出矛盾）$\alpha$ 是一条发散 WKB 曲线。文献 \cite{MR743423} 证明这样的曲线其实是回复的：令 $A$ 为 $\alpha$ 的闭包，则 $\alpha$ 无穷多次任意接近 $A$ 的每一点。此外，$A$ 的内部是一个非空连通区域，而 $A$ 的边界由有限 WKB 曲线组成。既然已约定不存在有限 WKB 曲线，唯一可能是 $A$ 的边界为空：换言之，$A$ 充满 $C$。特别地，$\alpha$ 会任意接近那些奇点。但回忆第 \ref{sec:small-flat} 节：附近的 WKB 曲线是进入奇点的对数螺线（除非 $e^{-i\vartheta} m \in i\IR$，此时它们是闭合的 BPS 曲线，而这一情形我们已排除）。因此 $\alpha$ 不能离奇点太近，否则就会掉进去。这就给出了所要的矛盾。（注意：这一论证关键地依赖于我们假定了 $n_{\rm sing} \ge 1$，即 $C$ 上有奇点。若 $\lambda^2$ 是正则的，WKB 叶状结构可以复杂得多。\footnote{这一点暗示：当 $C$ 上没有缺陷算符时，BPS 态的谱在定性上有所不同。的确，若 $C$ 上没有奇点，我们预期存在任意长的闭合 WKB 曲线。相应地，当没有缺陷算符时，我们预期存在一个无限谱的 BPS 矢量多重态，其质量任意大。}）于是我们的叶状结构只由一般 WKB 曲线与分离 WKB 曲线组成。

一般 WKB 曲线组成单参数族；每一族扫出一个「胞腔」，其边界是若干分离 WKB 曲线的并 \cite{MR743423}。在 $w$ 坐标下，这样一个胞腔看起来就是一条倾角为 $\vartheta$ 的带。再次利用不存在有限 WKB 曲线的假定：这条带的两条边界各自至多含一个转向点。微妙之处在于，这两个转向点可能实际上是同一个点。由此，胞腔的闭包只有两种可能的拓扑，如图 24 所示。% TODO 合稿改 \ref{fig:diamond-disc}
\footnote{若只有 $n_{\rm sing}=2$ 个奇点且 $C = \IC\IP^1$，则还会有第三种可能的胞腔，即整个 $C$；但我们只考虑 $n_{\rm sing} \ge 3$，从而排除了这一情形。}

\insfig{4in}{diamond-disc}{一般 WKB 曲线扫出的两种可能的胞腔：菱形与圆盘。圆盘可以看作退化的菱形，由把两个转向点等同起来而得到。一般 WKB 曲线用细黑线表示，分界线用粗红线表示。}

\subsection{定义 WKB 三角剖分}\label{subsec:DefiningWKB-Triang}

现在来定义 WKB 三角剖分。它大致与我们刚看到的胞腔分解「对偶」。

我们假定所有转向点都是 $\lambda^2$ 的单零点。在每一族中选一条一般 WKB 曲线 $E_i$。这些 $E_i$ 把 $C$ 分成一些面。每个面必含一个转向点（这可由下面的事实推出：这些边属于不同的胞腔，而每个胞腔的边界由分离 WKB 曲线组成）。于是让我们把注意力集中在单个转向点 $\CT$ 上。在 $\CT$ 附近跑的 WKB 曲线是某个面 $F$ 的边。看图 23 可知，% TODO 合稿改 \ref{fig:foliation-zero}
$\CT$ 至多在三个胞腔的边界上。若恰好是三个，则 $F$ 是三角形。也可能出现这样的情况：图 23 所画三类曲线中有两类其实属于同一个胞腔。% TODO 合稿改 \ref{fig:foliation-zero}
此时 $F$ 是退化三角形。

\insfig{4.5in}{optimal-faces}{WKB 三角剖分中出现的两类面。左边是真正的三角形。右边是退化三角形，只有两条不同的边和两个顶点（与图 19 比较）。% TODO 合稿改 \ref{fig:degenerate-face}
每个面含一个转向点。WKB 三角剖分的边用粗黑线表示，分界线用细红线表示。}

我们归纳出：$E_i$ 构成 $C$ 的一个三角剖分的边，其面只有两种类型，如图 25 所示。% TODO 合稿改 \ref{fig:optimal-faces}
该三角剖分的拓扑不依赖于我们在胞腔内选哪一条一般的 $E_i$，因此我们把仅差这种选择的两个三角剖分等同起来。也就是说，所谓「三角剖分」，我们真正指的是三角剖分的同痕类。在这个意义下，我们就定义了 $C$ 的一个三角剖分。

\subsection{WKB 装饰：小平坦截面} \label{sec:small-flat}

至此我们只用 WKB 叶状结构定义了一个未装饰的三角剖分。为了构造 $X^\vartheta_\gamma$，我们还需要某种典范的方式来选取装饰。

先考虑平坦截面在某个奇点附近的行为，为方便起见把该点放在 $z=0$。该点附近联络的行为由 (4.9) 给出，% TODO 合稿改 \eqref{eq:A-sing-reg}
它由 $\varphi$ 与 $A$ 的留数 $\rho$、$\alpha$ 分别决定。我们取 $\rho = 2m\sigma^3$，其中 $m \in \IC$；并取 $\alpha = -2i m^{(3)}\sigma^3$，其中 $m^{(3)} \in \IR$。

于是，按正则奇点附近行为的标准 Frobenius 分析（这里略有修改，因为我们考虑的是 $C^\infty$ 联络而非全纯联络），存在两个形如
\begin{equation} \label{eq:small-flat-1}
s^{(1)} = z^{-R \zeta^{-1} m + m^{(3)}} \bar{z}^{-R \zeta \bar{m} - m^{(3)}} \begin{pmatrix} 1 + O(\abs{z}) \\ O(\abs{z}) \end{pmatrix},
\end{equation}
\begin{equation} \label{eq:small-flat-2}
s^{(2)} = z^{R \zeta^{-1} m - m^{(3)}} \bar{z}^{R \zeta \bar{m} + m^{(3)}} \begin{pmatrix} O(\abs{z}) \\ 1 + O(\abs{z}) \end{pmatrix}
\end{equation}
的平坦截面。它们具有顺时针单值本征值 $\mu^{(1)} = e^{2\pi i\nu}$、$\mu^{(2)} = e^{-2\pi i\nu}$，其中
\begin{equation} \label{eq:nu}
\nu = R \zeta^{-1} m - 2 m^{(3)} - R \zeta \bar{m}.
\end{equation}
我们来计算它们沿一条进入奇点的 WKB 曲线的行为。利用 \eqref{eq:log-spiral} 得到
\begin{equation} \label{eq:small-flat-3}
s^{(1)} \sim \exp \left[ \varepsilon(\vartheta) \left\{ R (e^{i\vartheta}\zeta^{-1} + e^{-i\vartheta}\zeta) + m^{(3)} \left( \frac{e^{-i\vartheta}}{\bar m} - \frac{e^{i\vartheta}}{m} \right) \right\} t \right]
\end{equation}
其中 $\varepsilon(\vartheta)$ 已在 \eqref{eq:defeps} 中定义。乘 $m^{(3)}$ 的那一项是纯相位，不影响 $s^{(1)}$ 的范数。看剩下那一项可知：当 $\varepsilon(\vartheta)\,\Re (e^{i\vartheta}\zeta^{-1} + e^{-i\vartheta}\zeta) < 0$ 时，$\norm{s^{(1)}}$ 当 $t \to \infty$ 指数式小。类似地，当 $\varepsilon(\vartheta)\,\Re (e^{i\vartheta}\zeta^{-1} + e^{-i\vartheta}\zeta) > 0$ 时，$\norm{s^{(2)}}$ 当 $t \to \infty$ 指数式小。

因此，若 $\zeta \in \IH_\vartheta$，则顺时针单值为
\begin{equation} \label{eq:monodromy-formula}
\mu = e^{(-\varepsilon(\vartheta)) 2\pi i\nu} = \exp \left[(-\varepsilon(\vartheta)) 2\pi i \left( R \zeta^{-1} m - 2 m^{(3)} - R \zeta \bar{m} \right) \right]
\end{equation}
的那个本征截面，沿进入奇点的 WKB 曲线其范数渐近地小于所有其他平坦截面。我们称它为「小平坦截面」。第 7.4 节将说明：% TODO 合稿改 \ref{sec:asymptotics}
若 $\zeta \to 0$ 且始终留在 $\IH_\vartheta$ 内，WKB 近似能很好地控制这个截面沿 WKB 曲线的演化。出于这个动机，我们选这个小平坦截面作为奇点处的典范装饰。

至此我们就对每个 $(\vartheta, \lambda^2)$ 定义好了一个装饰三角剖分 $T_\WKB(\vartheta, \lambda^2)$。我们称它为\emph{WKB 三角剖分}。为记号简便，有时把它写成 $T_\WKB(\vartheta)$，甚至只写 $T_\WKB$。回忆二次微分 $\lambda^2$ 对应于点 $u \in \CB$，所以我们有时也写 $T_\WKB(\vartheta, u)$。

\subsection{WKB 三角剖分的跳变}\label{subsec:Jumps-WKB-triang}

当我们变动参数 $\vartheta$ 时，WKB 叶状结构随之改变，相应地 $T_\WKB(\vartheta)$ 也改变。对一般的 $\vartheta$，$T_\WKB(\vartheta)$ 只是各条边作连续同伦，因而三角剖分的同痕类不变。但也有某些特殊值 $\vartheta = \vartheta_c$，在那里 $T_\WKB(\vartheta)$ 发生不连续的跳变。

这些跳变对我们至关重要：它们与 BPS 态的存在有关。的确，$T_\WKB$ 的拓扑完全由分离 WKB 曲线的行为决定，而这一行为恰好在一条分界线退化到含有一段有限 WKB 曲线时发生跳变。\footnote{考虑一条从 $z_0$ 出发的分离 WKB 曲线。在 $w$-平面中（$w = \int_{z_0}^z \lambda$ 在该分离曲线的某个邻域内有定义），这条曲线是一条直线。它随 $\vartheta$ 连续变化。这一点只会在如下情形失效：该分离 WKB 曲线的任何开邻域都含有另一个转向点。也就是说，它只在分离 WKB 曲线包含一段有限 WKB 曲线时失效。}现在把 (3.11) 与 \eqref{eq:wkb-curve} 作比较，% TODO 合稿改 \eqref{eq:BPS-WKB-curve}
可知 BPS 弦与 WKB 曲线服从完全相同的方程；有限 WKB 曲线正是总质量有限的 BPS 弦。于是，WKB 叶状结构改变拓扑的那些值 $\vartheta_c$ 就是 BPS 态的相位。

下面我们逐一考察当 $\vartheta$ 变化时 $T_\WKB(\vartheta)$ 可能发生的三种拓扑变化。

\subsubsection{来自 BPS 超多重态的跳变} \label{sec:hyper-jump}

最基本的拓扑变化例子发生在这样的 $\vartheta = \vartheta_c$ 处：此时出现一条连接两个转向点的 WKB 曲线。如第 3.1.3 节所述，% TODO 合稿改 \ref{sec:bps-strings}
这条有限 WKB 曲线代表 $d=4$ 理论的一个 BPS 超多重态。

当 $\vartheta$ 越过 $\vartheta_c$ 时，WKB 叶状结构经历一次拓扑变化；见图 26。% TODO 合稿改 \ref{fig:hypermultiplet-flip}
特别地，那条从西北走向东南的一般 WKB 曲线被一条从东北走向西南的曲线取代。既然这条一般 WKB 曲线代表 $T_\WKB$ 的一条边 $E$，我们就看到：当 $\vartheta$ 越过 $\vartheta_c$ 时，这个三角剖分经历一次 \emph{flip}。
\insfig{6in}{hypermultiplet-flip}{当 $\vartheta$ 越过某个临界值 $\vartheta_c$（此时出现一条有限 WKB 曲线，对应一个 BPS 超多重态）时，WKB 叶状结构的跳变。}

\subsubsection{$m e^{-i\vartheta} \in i\IR$ 时的跳变} \label{sec:flavor-jump}

更复杂的跳变行为发生在存在闭合 WKB 曲线的那些相位 $\vartheta = \vartheta_c$ 处。如第 3.1.3 节所述，% TODO 合稿改 \ref{sec:bps-strings}
这些闭合 WKB 曲线总是以单参数族出现。

最简单的情形是环绕单个奇点的一族闭合 WKB 曲线。正如第 \ref{sec:wkb-local} 节已看到的，只要 $m e^{-i\vartheta} \in i\IR$，就会出现这样一族。WKB 叶状结构在 $\vartheta$ 取该值附近的行为如图 27 所示。% TODO 合稿改 \ref{fig:spiral-jump}

\insfig{5.6in}{spiral-jump}{当 $\vartheta$ 越过某个临界值 $\vartheta_c$（此时出现一族环绕单个奇点的闭合 WKB 曲线）时，WKB 叶状结构的跳变。}

第 \ref{sec:small-flat} 节中我们定义了 WKB 三角剖分的装饰，它由 $\Re (e^{-i\vartheta} m)$ 的符号 $\varepsilon(\vartheta)$ 决定。这个符号在 $\vartheta_c$ 处改变。由此可知，在 $\vartheta_c$ 处 WKB 三角剖分经历一次 \emph{pop}。

\subsubsection{来自 BPS 矢量多重态的跳变} \label{sec:vector-jump}

另一种可能是：一族闭合 WKB 曲线并不收缩到单个奇点上。这对应于带某个非零规范荷的 BPS 矢量多重态。

一般地，这样的矢量多重态出现时，伴随有两族无穷多的超多重态：若取一个包含 $\vartheta_c$ 的足够窄的区间 $\vartheta_- < \vartheta < \vartheta_+$，则 WKB 三角剖分随 $\vartheta$ 光滑变化，唯独在 $C$ 的一个环形区域里例外，而在这个环形区域内 WKB 三角剖分经历无穷多次 flip。这正是第 5.9 节所考虑的情形。% TODO 合稿改 \ref{sec:limit-triangulation}
从 $\vartheta_+$ 出发让 $\vartheta$ 减小，$T_\WKB(\vartheta, \lambda^2)$ 依次经过一列三角剖分 $T_m$（$m>0$），且当 $\vartheta \to \vartheta_c^+$ 时 $m \to \infty$。另一方面，若改从 $\vartheta_-$ 出发，则当 $\vartheta$ 增大时 $T_\WKB(\vartheta, \lambda^2)$ 依次经过所有 $m<0$ 的 $T_m$，且当 $\vartheta \to \vartheta_c^-$ 时 $m \to -\infty$。

\insfig{5.4in}{vectormultiplet-twist}{WKB 叶状结构的一个环形区域，靠近某个临界值 $\vartheta = \vartheta_c$——此时出现一族代表 BPS 矢量多重态的闭合 WKB 曲线。}

大致说来，经过 $m$ 次 flip 之后，一条典型的 WKB 曲线会绕环带跑 $m$ 圈，方向是顺时针还是逆时针取决于 $\vartheta > \vartheta_c$ 还是 $\vartheta < \vartheta_c$。恰在 $\vartheta = \vartheta_c$ 处，WKB 曲线变成布满环带内某个开区域的闭合曲线。见图 28。% TODO 合稿改 \ref{fig:vectormultiplet-twist}

\subsubsection{其他退化是非一般的}

现在简略论证我们的断言：对一般的二次微分 $\lambda^2$，当 $\vartheta$ 变化时 WKB 三角剖分只会出现上述三种退化。（对非一般的模，$T_\WKB(\vartheta, \lambda^2)$ 的其他退化可以并且确实会发生。下面第 10.2 节会看到一个例子。）% TODO 合稿改 \ref{sec:nf1}

一个临界值 $\vartheta_c$ 出现在一条分离 WKB 曲线退化到含有一段有限 WKB 曲线之时，而这段有限 WKB 曲线起止于一个转向点。只有两种可能：这两个转向点相同，或者不同。若两个转向点不同，则这条有限 WKB 曲线没有模，$\vartheta_c$ 附近的行为必是第 \ref{sec:hyper-jump} 节的那种超多重态退化。若两个转向点相同，则通过无穷小位移点（位于叶状结构中局部以此有限 WKB 曲线为边界的那一胞腔内）的 WKB 曲线必定闭合。这些闭合 WKB 曲线构成单参数族。这一族可能在以下情形终止：闭合 WKB 曲线碰到奇点、收缩至零尺寸、或碰到转向点。考察奇点附近的局部行为（对一般的 $m$）可知，闭合 WKB 曲线族不可能以碰到奇点的方式终止。若闭合 WKB 曲线收缩至零尺寸，则我们处于第 \ref{sec:flavor-jump} 节的情形。最后一种可能是闭合 WKB 曲线碰到另一个转向点。一般情形下它会恰好碰到一个转向点。于是我们的族所描述的，正是第 \ref{sec:vector-jump} 节退化中出现的情形。
